\subsubsection{Término de costura en la concatenación \(108\to1080\)}
\label{vi:subsubsec:costura-108-1080}

La prolongación de la lectura cronológica debe conservar el predecesor
efectivo de cada evento. Sean
\(d^{(b)}=(d^{(b)}_1,\ldots,d^{(b)}_{108})\), \(0\le b<10\), diez
palabras consecutivas, y definamos
\[
 a_b=d^{(b)}_1,\qquad q_b=d^{(b)}_{108},\qquad q_{-1}=q_9.
\]
Denotemos por \(T_z^{\sqcup}\) la suma de los diez lectores evaluados
por separado, usando \(q_b\) como predecesor cíclico de \(a_b\), y por
\(T_z^{\Vert}\) el lector de la palabra concatenada, que usa
\(q_{b-1}\) como predecesor efectivo de \(a_b\).

\begin{lemma}[Corrección exacta de costura]
\label{vi:lem:costura-108-1080}
Con \(\Delta T_z=T_z^{\Vert}-T_z^{\sqcup}\),
\begin{equation}
 \Delta T_z=
 \sum_{b=0}^{9}\mathbf1_{\{a_b=9\}}
 \bigl(j_2(q_{b-1})-j_2(q_b)\bigr).
 \label{vi:eq:delta-costura-tz}
\end{equation}
Si las coronas se sitúan en \(t\equiv0\pmod{27}\), entonces
\begin{equation}
 \Delta C_z=0,
 \qquad
 \Delta k_{\Delta_4}=\Delta T_z.
 \label{vi:eq:delta-costura-k}
\end{equation}
\end{lemma}

\begin{proof}
Fuera de los diez comienzos de bloque, cada evento conserva el mismo
predecesor en ambas evaluaciones. Si \(a_b\ne9\), el indicador que define
\(z_t\) anula también la posible diferencia de predecesor. Si \(a_b=9\),
el término separado es \(j_2(q_b)\) y el concatenado es
\(j_2(q_{b-1})\); su diferencia es el sumando de
\eqref{vi:eq:delta-costura-tz}. Al sumar los diez bloques se obtiene la
primera identidad. Los nuevos comienzos ocupan las posiciones
\(108b+1\equiv1\pmod{27}\), nunca una corona. Por tanto la costura no
modifica \(C_z\), y la definición
\(k_{\Delta_4}=T_z-2C_z\) da
\eqref{vi:eq:delta-costura-k}.
\end{proof}

La corrección alcanza realmente el valor \(-2\). Basta tomar un único
comienzo \(a_b=9\), con \(q_{b-1}=1\), \(q_b=2\), y exigir
\(a_r\ne9\) en los demás comienzos. Entonces
\(j_2(q_{b-1})-j_2(q_b)=-1-1=-2\). Esta palabra es un control del lector
de costura; no se promueve a ruta superviviente, partícula ni firma
física.

La longitud \(1080\) repite diez estaciones de fase, no diez copias de
una profundidad invariable. El transporte correcto es la composición
ordenada
\[
 \mathcal H_9\circ\mathcal H_8\circ\cdots\circ\mathcal H_0.
\]
Sólo una prueba adicional de estacionariedad permitiría sustituirla por
\(\mathcal H_{108}^{10}\). La fórmula de costura conserva precisamente
la memoria que esa sustitución podría borrar.
